\documentclass[11pt]{article}

% --- Packages and formatting ---
\usepackage[margin=1in]{geometry}
\usepackage[T1]{fontenc}
\usepackage{lmodern}
\usepackage{microtype}
\usepackage{graphicx}
\usepackage{xcolor}
\usepackage{hyperref}
\usepackage{enumitem}
\usepackage{parskip}
\hypersetup{colorlinks=true,linkcolor=blue,urlcolor=blue,citecolor=blue}
\definecolor{promptgray}{HTML}{666666}

% Replace each \prompt{...} with your own prose before submitting.
\newcommand{\prompt}[1]{\textcolor{promptgray}{\textit{[#1]}}}

\title{\textbf{An Informative Project Title}}
\author{Student 1 Name \quad \& \quad Student 2 Name \\
        CSC 245/445 --- Team Name}
\date{}

\begin{document}
\maketitle
\vspace{-1.5em}

\section{Problem \& Motivation}

\subsection{Problem Statement and Motivation}
\prompt{Clearly describe the problem you want to solve, why it is important,
who would benefit, and what motivates your choice. You may draw inspiration
from open challenges and workshops at ICLR, NeurIPS, CVPR, or ACL, or from
applications of deep learning that interest you.}

\subsection{ML/DL Task Formulation}
\prompt{Explain how you formulate the problem as a machine-learning or
deep-learning task. Specify the input, desired output, and learning setting
(e.g., classification, regression, generation, retrieval, segmentation).}

\subsection{Illustrative Example (Optional)}
\prompt{Add a small illustrative figure, example input/output pair, or task
diagram if it helps the reader understand the problem.}

% To insert a figure, uncomment and edit:
% \begin{figure}[htbp]
%   \centering
%   \includegraphics[width=0.75\linewidth]{problem-illustration.pdf}
%   \caption{Illustration of the proposed task.}
%   \label{fig:problem}
% \end{figure}

\section{Methods}

\subsection{Open-Source Baseline}
\prompt{Identify an existing open-source model or method expected to
perform reasonably well on your task. Cite the paper and link the
implementation, and explain why it is an appropriate baseline.}

\subsection{Proposed Method}
\prompt{Describe your proposed approach and the concrete steps needed to
implement it. If the problem is new, identify an existing method that can
reasonably be adapted or extended to solve it. Explain how your approach
relates to the baseline.}

\subsection{Research Hypothesis (CSC 445)}
\prompt{For CSC 445: state a specific, testable research hypothesis involving
an innovative improvement over the baseline. Explain why you expect it to
work. For CSC 245, omit this subsection if it is not required.}

\subsection{Method Illustration (Optional)}
\prompt{Include a baseline architecture diagram or method figure if it
would clarify the approach.}

% To insert a baseline/model figure, uncomment and edit:
% \begin{figure}[htbp]
%   \centering
%   \includegraphics[width=0.8\linewidth]{baseline-architecture.pdf}
%   \caption{Baseline model or proposed method overview.}
%   \label{fig:method}
% \end{figure}

\section{Evaluation Plan}

\subsection{Datasets and Test Cases}
\prompt{Identify the dataset(s) on which you will evaluate your method.
Explain the data source, approximate size, train/validation/test split,
and relevant preprocessing. Consider challenging or unusual corner cases
that may be overlooked in existing literature.}

\subsection{Evaluation Metrics}
\prompt{Specify the performance measure(s) you will report and justify
why each metric is appropriate for the task (e.g., accuracy, F1, AUROC,
BLEU, IoU, MAE).}

\subsection{Training / Fine-Tuning Feasibility}
\prompt{Provide a realistic plan for model training. For smaller models,
explain how you will train on the dataset. For large/foundation models
(e.g., LLMs or ViTs), explain how you will fine-tune them. Include
compute resources, time constraints, and a feasible experimental scope.}

\subsection{Comparison Strategy}
\prompt{Explain how you will compare the baseline with your proposed
method, and what results would support or reject your hypothesis.}

% Add bibliography entries or your preferred BibTeX setup when needed.
% \bibliographystyle{plain}
% \bibliography{references}

\end{document}